% !TeX root = closed-systems-from-comparison-completeness.tex
% ============================================================
\chapter*{Orientation to Part V}
\phantomsection
\label{part:intro-curvature}
% ============================================================
\begin{chapterorientation}
\orientationitem{Settled}{The transport obstruction and stitched arena \(ST\) have been constructed internally.}
\orientationitem{Task}{Part~V follows the obstruction into curvature, causal scaling, Hilbert structure, macroscopic compatibility, and connection-first geometry.}
\orientationitem{Handoff}{The output is the realized geometric and Hilbert-side structure needed for phase curvature.}
\end{chapterorientation}

Part~V is the main bridge from intrinsic transport to realized physical
geometry.  It does not assume a background manifold as primitive.  The
smooth language enters only as a realization of transport and
comparison-side structure already obtained.

\Cref{ch:christoffel,ch:interface} identify relational loop defect with
curvature and stabilize the quadratic carrier on the Stone-limit side.
\Cref{ch:causality,ch:hilbert} derive the causal-scaling and Hilbert
structures carried by that quadratic layer.  \Cref{ch:einstein}
states the macroscopic compatibility law in the realized branch, and
\cref{ch:connection} reconstructs geometry from connection data.  The
part ends with the geometric and phase-ready structure required for the
electromagnetic sector.
